\chapter{使用说明、地基与索引}

\section{这本手册解决什么问题}

这不是算法百科，而是区域赛现场的“关键词到代码”工具。看到题面后，先判断对象和约束，再翻到对应卡片。每张卡片都从上级目录的 \texttt{muban.cpp} 复制，随后追加该算法自己的结构和 \texttt{solve()} 接口。

\begin{card}
\textbf{现场流程：}
\begin{enumerate}
\item 读清输入规模、是否多测、下标、取模和无解要求；
\item 用本章索引定位算法；
\item 复制 \texttt{muban.cpp}，再复制算法卡片中的追加代码；
\item 按“接入提示”改变量含义；
\item 先跑卡片样例，再跑边界测试。
\end{enumerate}
\end{card}

\section{统一代码地基}

统一地基就是 \texttt{muban.cpp}，不再维护第二套头文件。当前版本的约定是：

\begin{lstlisting}[style=cpp]
#include <bits/stdc++.h>
using namespace std;
#define int long long
#define endl '\n'
using i64 = long long;
using i128 = __int128_t;
// 别把这个改成具体的数,比赛的时候谁有时间给你重写
const int inf = 2e18+10;
const int N = 1e6 + 10;
const double eps = 1e-8;
\end{lstlisting}

本手册默认使用 \texttt{int}（宏展开为 \texttt{long long}），不处理需要 \texttt{unsigned} 的模板。算法卡片必须说明自己是否需要覆盖 \texttt{N}、\texttt{inf}、\texttt{eps}，以及新增的节点池大小。

\begin{warnbox}
\texttt{\#define int long long} 后主函数写 \texttt{signed main()}；位移写 \texttt{1LL << k}。数组下标仍然要检查范围，宏不会替代边界检查。
\end{warnbox}

\section{复杂度速查}

\begin{center}
\begin{tabular}{p{0.18\linewidth}|p{0.27\linewidth}|p{0.42\linewidth}}
\textbf{规模} & \textbf{常见可接受复杂度} & \textbf{优先考虑}\\ \hline
\(n\le 20\) & \(O(2^n n)\) & 状压、子集枚举\\
\(n\le 40\) & \(O(2^{n/2})\) & 折半搜索\\
\(n\le 200\) & \(O(n^3)\) & Floyd、区间 DP、矩阵\\
\(n\le 2\times 10^5\) & \(O(n\log n)\) & 排序、树状数组、线段树、Dijkstra\\
\(n\le 10^6\) & \(O(n)\) 或 \(O(n\log\log n)\) & 扫描、哈希、线性筛\\
\(n\) 很大 & \(O(\log n)\) 或 \(O(\log^2 n)\) & 快速幂、矩阵、数论分块
\end{tabular}
\end{center}

\section{关键词索引}

\begin{tabular}{p{0.38\linewidth}|p{0.42\linewidth}}
\textbf{题面信号} & \textbf{先查}\\ \hline
区间加、区间和 & 懒标记线段树\\
单点加、前缀和、第 \(k\) 小 & 树状数组 / 可持久化线段树\\
静态区间最值 & 稀疏表\\
离线区间统计 & 莫队\\
历史版本、版本差异 & 可持久化结构\\
无权最短路 & BFS / 多源 BFS\\
边权只有 \(0,1\) & 0-1 BFS\\
非负权最短路 & Dijkstra\\
负边或负环 & Bellman--Ford / 差分约束\\
有向图缩点 & SCC\\
真假约束、或关系 & 2-SAT\\
多模式串 & AC 自动机\\
不同子串、出现次数 & SAM\\
后缀排序、LCP & 后缀数组\\
树上路径 & LCA / HLD\\
子树颜色统计 & DSU on Tree\\
树上距离在线查询 & 点分治 / 点分树\\
\(x^n\)、线性递推 & 快速幂 / 矩阵快速幂\\
多个同余方程 & CRT / 扩展 CRT\\
凸包、最远点对 & Andrew / 旋转卡壳\\
矩形并面积 & 扫描线 + 线段树
\end{tabular}

\section{卡片格式}

每个算法在两本手册中都按以下顺序出现：

\begin{enumerate}
\item 识别关键词；
\item 适用条件和不能使用的情况；
\item 复杂度、内存、下标约定；
\item 状态或节点含义；
\item 追加代码；
\item 从 \texttt{muban.cpp} 接入的方式；
\item 完整样例输入输出；
\item 四组最小边界测试；
\item 常见错误和改造点。
\end{enumerate}
